\begin{ouverture}[calcul]{7}{Utiliser et comparer les fractions}{Partie 2 · Fractions et nombres décimaux}
{Demi-droite · Comparaisons · Sommes de parts}
\objectif{Placer et repérer des fractions sur une demi-droite graduée.}
\objectif{Comparer des fractions et compléter jusqu’à l’unité.}
\objectif{Additionner et soustraire des parts de même taille.}
\prerequis{Fractions, numérateur et dénominateur.}
\end{ouverture}
\pagestyle{otMain}\essentiel\fontsize{11}{15}\selectfont
\phantomsection\label{cours-7-1}
\essSection{1}{Des fractions sur la demi-droite}
\begin{propriete}Sur une demi-droite graduée, les fractions sont des nombres, placés entre les entiers. Si l’unité est partagée en 4, chaque pas vaut $\frac14$ : le point $\frac74$ se trouve 3 pas après 1.\end{propriete}
\label{fig-cours-7-1-cours-droite-fractions}\figureGeo{cours-droite-fractions}

\begin{savoirfaire}[Mesurer avec des fractions]Pour mesurer avec une bande unité : reporter la bande entière, puis utiliser ses plis en quarts pour finir. Une longueur de 2 unités et 3 quarts s’écrit $2+\frac34$, c’est-à-dire $\frac{11}4$ d’unité.\end{savoirfaire}
\begin{attention}Toujours lire le partage de l’unité avant de placer une fraction : $\frac32$ n’est pas au même endroit si l’unité est partagée en 2 ou en 4. Entre 0 et 1, il y a une infinité de fractions.\end{attention}
\clearpage
\phantomsection\label{cours-7-2}
\essSection{2}{Comparer et réunir des parts}
\begin{propriete}Avec le même dénominateur, les parts ont la même taille : la fraction la plus grande a le plus grand numérateur, et on peut réunir ou retirer des parts : $\frac38+\frac28=\frac58$ et $\frac7{10}-\frac4{10}=\frac3{10}$.\end{propriete}
\begin{exemple}Avec le même numérateur, on compare la taille des parts : $\frac13>\frac14$, car partager en 3 donne de plus grandes parts que partager en 4. Donc $\frac23>\frac24$.\end{exemple}
\begin{savoirfaire}[Prendre une part d’une quantité]Prendre $\frac13$ de 12 billes, c’est partager 12 en 3 : $12:3=4$ billes. Un quart de 100 mètres vaut $100:4=25$ mètres. La fraction unitaire agit comme un partage.\end{savoirfaire}
\begin{savoirfaire}[Compléter jusqu’à 1]Il manque $\frac49$ pour aller de $\frac59$ à l’unité, car $\frac59+\frac49=\frac99=1$. Penser « combien de parts manquent ? ».\end{savoirfaire}

\clearpage\rubrique{Découverte}
\begin{blocLivre}{ex-7-1}\begin{exercice}[1]{Lire sur la demi-droite}
L’unité est partagée en quarts. Lire les fractions repérées par les points A et B.
\label{fig-ex-7-1}\figureGeo{fd-1}
\end{exercice}\end{blocLivre}
\begin{blocLivre}{ex-7-2}\begin{exercice}[1]{Placer des tiers}
Sur la demi-droite graduée en tiers, placer les points $C\left(\frac23\right)$ et $D\left(\frac53\right)$.
\label{fig-ex-7-2}\figureGeo{fd-2}
\end{exercice}\end{blocLivre}
\begin{blocLivre}{ex-7-3}\begin{exercice}[1]{Comparer à l’unité}
Comparer chaque fraction à 1 : $\frac56$ ; $\frac76$ ; $\frac66$.
\end{exercice}\end{blocLivre}
\begin{blocLivre}{ex-7-4}\begin{exercice}[1]{Même dénominateur}
Comparer $\frac58$ et $\frac38$, puis ranger dans l’ordre croissant : $\frac45$ ; $\frac15$ ; $\frac75$.
\end{exercice}\end{blocLivre}
\begin{blocLivre}{ex-7-5}\begin{exercice}[1]{Une part d’une quantité}
Calculer $\frac13$ de 12 billes, puis $\frac14$ de 100 mètres.
\end{exercice}\end{blocLivre}

\rubrique{Entraînement}
\begin{blocLivre}{ex-7-6}\begin{exercice}[2]{Réunir des parts}
Calculer $\frac38+\frac28$, puis $\frac25+\frac35$.
\end{exercice}\end{blocLivre}
\begin{blocLivre}{ex-7-7}\begin{exercice}[2]{Retirer des parts}
Calculer $\frac7{10}-\frac4{10}$, puis $1-\frac29$.
\end{exercice}\end{blocLivre}
\begin{blocLivre}{ex-7-8}\begin{exercice}[2]{Compléter jusqu’à 1}
Compléter : $\frac59+\frac{\square}9=1$, puis $\frac{\square}6+\frac16=1$.
\end{exercice}\end{blocLivre}
\begin{blocLivre}{ex-7-9}\begin{exercice}[2]{Des cinquièmes bien placés}
L’unité est partagée en cinquièmes. Lire la fraction du point E, puis placer $F\left(1+\frac25\right)$.
\label{fig-ex-7-9}\figureGeo{fd-9}
\end{exercice}\end{blocLivre}
\begin{blocLivre}{ex-7-10}\begin{exercice}[2]{Quelle part est la plus grande ?}
Comparer $\frac13$ et $\frac14$ à l’aide des bandes, puis comparer $\frac23$ et $\frac24$.
\label{fig-ex-7-10}\figureGeo{fd-10}
\end{exercice}\end{blocLivre}

\rubrique{Problèmes}
\begin{blocLivre}{ex-7-11}\begin{exercice}[3]{Mesurer avec une bande}
Avec la bande unité pliée en quarts, mesurer la longueur du segment tracé.
\label{fig-ex-7-11}\figureGeo{fd-11}
\end{exercice}\end{blocLivre}
\begin{blocLivre}{ex-7-12}\begin{exercice}[3]{La randonnée}
Un sentier mesure 20 km. Le matin, le groupe parcourt le quart du sentier.\begin{questions}
\item Quelle distance est parcourue le matin ?
\item Quelle distance reste-t-il pour l’après-midi ?
\end{questions}
\end{exercice}\end{blocLivre}
\begin{blocLivre}{ex-7-13}\begin{exercice}[3]{Sept parts de deux tartes}
Deux tartes identiques sont coupées en sixièmes et 7 parts sont mangées. Quelle fraction de tarte reste-t-il ?
\end{exercice}\end{blocLivre}
\begin{blocLivre}{ex-7-14}\begin{exercice}[3]{La même longueur ?}
Lila plie une bande en 2 et prend 1 part ; Noa plie la même bande en 4 et prend 2 parts. Ont-ils pris la même longueur ? Justifier avec les bandes.
\label{fig-ex-7-14}\figureGeo{fd-14}
\end{exercice}\end{blocLivre}
\begin{blocLivre}{ex-7-15}\begin{exercice}[3]{Le ruban partagé}
Un ruban de 18 m est partagé : Ana reçoit la moitié, Bilal reçoit le tiers, Chloé reçoit le reste.\begin{questions}
\item Calculer les longueurs reçues par Ana et Bilal.
\item Quelle longueur reste-t-il pour Chloé ?
\end{questions}
\end{exercice}\end{blocLivre}

\clearpage
